\newif\ifglobal

%%%%%%%%%%%%%%%%%%%%%%%
%%% Compile to view %%%
%%%%%%%%%%%%%%%%%%%%%%%

%%%%%%%%%%%%%%%%%%%%%%%%%%%%%%%%%%%%%%%%%%%%%%%%%%%
%%% Readme:                                     %%%
%%% https://www.overleaf.com/read/bwdjvfjksvjy/ %%%
%%%%%%%%%%%%%%%%%%%%%%%%%%%%%%%%%%%%%%%%%%%%%%%%%%%

% >>> M?: marks a module setting number or important notice

% >>> M4: Set {./relative-path} to external files for this module <<<
\def \modulefiles {./30-RNG}
% To add an external file, use:
% (Replace '---' with actual filename)
% '\input{\modulefiles/tables/---.tex}' for tables
% '\includegraphics{\modulefiles/figures/---}' for figures
% '\subfile{\modulefiles/---}' for register files

% >>> M1: This module is in global project? <<<
% 'YES' - uncomment; 'NO' - comment out
\globaltrue

\ifglobal
    \documentclass[main\_\_EN.tex]{subfiles}
   
\else
    \documentclass[openany, 10pt]{book}

    % >>> M2: In ./00-shared/preamble-trm-module, also find
    %         settings for visibility of todo notes, watermark <<<
    \usepackage{preamble-trm-module}

    % >>> M3: Temporary labels in English,
    %         you can add your own labels there <<<
    \myexternaldocument{./00-shared/temp-labels-en}

\fi %\ifglobal


\setcounter{secnumdepth}{4}


% Define formatting of parts and chapters in text
\titleformat{\chapter}
  [display]
  {\LARGE\bfseries}
  {\color{AllportsColor}Chapter \thechapter}
  {1em}
  {\color{AllportsColor}}
\titlespacing*{\chapter}{0pt}{0pt}{20pt}

% Define headers
\usepackage{fancyhdr}
\usepackage{lipsum}
\renewcommand{\chaptermark}[1]{\markboth{Chapter\, \thechapter\, \ #1}{}}
\fancypagestyle{plain}{%
  \fancyhead{}%
  %\fancyhead[R]{\Navigation}
  \fancyhead[L]{%
    \nouppercase{%
      \ %
      \normalsize\leftmark
    }%
  }%
}
\pagestyle{plain}


% >>> M5: If this module needs any updates in
% './00-shared/preamble-trm-module.sty'
% (include additional package, etc.), contact your module PM
% or create the file './00-shared/changelog.tex' make the
% necessary updates yourself and document them in this file <<<

% Make text left-aligned and allow latex to break pages naturally, comment out for full justification
\raggedright
\raggedbottom


\begin{document}


% Sets language into which pre-defined bits of text are translated
\selectlanguage{english}

% Do not delete this block
\ifglobal
% Add the following front matter if
% this module is outside of global project
\else
    \InputIfFileExists{readme}{}{}
    \listoftodos
    {\let\clearpage\relax \listofdonetodos}
    \tableofcontents
    %\listoftables
    %\listoffigures
\fi


%%%%%%%%%%%%%%%%%%%%%%%%%
%%% TRM Module Begins %%%
%%%%%%%%%%%%%%%%%%%%%%%%%

\hypertarget{rng}{}
\chapter{Random Number Generator (RNG)}\label{mod:rng}

\section{Introduction}

The \chipname{} contains a true random number generator, which generates 32-bit random numbers that can be used for cryptographical operations, among other things.

\section{Features}

The random number generator in \chipname{} generates true random numbers, which means random numbers derived from a physical process, rather than by means of an algorithm. All generated random numbers have an equal probability of appearing within a specific range.

\section{Functional Description}\label{sec:rng-func-descr}

Every 32-bit value read from the \hyperref[regdesc:LPPERIRNGSYNCDATAREG]{LPPERI\_RNG\_DATA\_SYNC\_REG} register of the random number generator is a true random number. These true random numbers are generated based on the \textbf{thermal noise} in the system, the \textbf{asynchronous clock mismatch}, and the \textbf{asynchronous counter}.

\begin{itemize}
    \item \textbf{Thermal noise} comes from the high-speed ADC, or SAR ADC, or both. Whenever the high-speed ADC or SAR ADC is enabled, bit streams will be generated and fed into the random number generator through an XOR logic gate as random seeds.
    \item RC\_FAST\_CLK is an \textbf{asynchronous clock source} that can generate metastability in the circuit. This metastability can be used as a random number entropy to feed into the random number generator.
    \item \textbf{Asynchronous counters} count using an asynchronous clock source. The asynchronous clock source generates circuit metastability, which can be used as a random number seed to feed into the random number generator. There are two types of asynchronous counters:
    \begin{itemize}
        \item On-chip \textbf{\nameref{mod:rtctimer} (RTC\_TIMER)}: The counter directly counts using the on-chip RTC timer. The count value is fed into the random number generator as a random number seed after applying XOR logic.
        \item \textbf{Ring buffer (BUF\_CHAIN)} as the noise source.
    \end{itemize}
    These two types of asynchronous counters can work independently or in combination.
\end{itemize}

\begin{tiplisting}
    For limitations of particular noise sources, please check \hyperref[others:rng-note-clock-seed]{notes in Section Programming Procedures}.
\end{tiplisting}

\begin{figure}[ht]
    \centering
    \includegraphics[width=0.8\textwidth]{\modulefiles/figures/29-RNG-noise-source}
    \caption{Noise Source}
    \label{fig:rng-noise-source}
\end{figure}

When there is noise coming from the SAR ADC, the random number generator is fed with a 2-bit entropy in one clock cycle of RC\_FAST\_CLK, which is generated from an internal RC oscillator (see Chapter \ref{mod:resclk} \textit{\nameref{mod:resclk}} for details). Thus, it is advisable to read the \hyperref[regdesc:LPPERIRNGSYNCDATAREG]{LPPERI\_RNG\_DATA\_SYNC\_REG} register at a maximum rate of 1 MHz to obtain the maximum entropy.

When there is noise coming from the high-speed ADC, the random number generator is fed with a 2-bit entropy in one APB clock cycle, which is normally 80 MHz. Thus, it is advisable to read the \hyperref[regdesc:LPPERIRNGSYNCDATAREG]{LPPERI\_RNG\_DATA\_SYNC\_REG}
register at a maximum rate of 5 MHz to obtain the maximum entropy.

A data sample of 2 GB, which is read from the random number generator at a rate of 5 MHz with only the high-speed ADC being enabled, has been tested using the Dieharder Random Number Testsuite (version 3.31.1). The sample passed all tests.


\section{Programming Procedure}

To use the random number generator, set \hyperref[fielddesc:LPPERIRNGCKEN]{LPPERI\_RNG\_CK\_EN} to enable its clock. The clock is automatically enabled when \nameref{mod:ecdsa} is used, regardless of 
\hyperref[fielddesc:LPPERIRNGCKEN]{LPPERI\_RNG\_CK\_EN}.

When using the random number generator, ensure the entropy source is enabled. Otherwise, pseudo-random numbers will be returned. 
\begin{itemize}
    \item SAR ADC can be enabled by using the DIG ADC controller. For details, please refer to Chapter \ref{mod:adc} \textit{\nameref{mod:adc}}.
    \item High-speed ADC is enabled automatically when the Wi-Fi or Bluetooth module is enabled.
    \item RC\_FAST\_CLK is always enabled
    \item Asynchronous counters
        \begin{itemize}
            \item RTC timer (RTC\_TIMER) is enabled by setting the \hyperref[fielddesc:LPPERIRTCTIMEREN]{LPPERI\_RTC\_TIMER\_EN} field.
            \item Ring buffer (BUF\_CHAIN) is enabled by setting the \hyperref[fielddesc:LPPERIRNGSAMPLEENABLE]{LPPERI\_RNG\_SAMPLE\_ENABLE} bit.
        \end{itemize}
\end{itemize}

\begin{tiplisting}\label{others:rng-note-clock-seed}
\vspace{-2.5em}
\begin{enumerate}
    \item When the Wi-Fi module is enabled, the value read from the high-speed ADC can be saturated in some extreme cases, which lowers the entropy. Thus, it is advisable to also enable the SAR ADC as the noise source for the random number generator for such cases.
    \item Enabling RC\_FAST\_CLK and asynchronous counters increases the RNG entropy. However, to ensure maximum entropy, it’s recommended to always enable the SAR ADC or high-speed ADC as well.
\end{enumerate}
\end{tiplisting}

Read the \hyperref[regdesc:LPPERIRNGSYNCDATAREG]{LPPERI\_RNG\_DATA\_SYNC\_REG} register multiple times until sufficient random numbers have been generated. Ensure the rate at which the register is read does not exceed the frequencies described in Section \ref{sec:rng-func-descr} above.


%%%%%%%%%%%%%%%%%%%%%%%%
%%% Register Summary %%%
%%%%%%%%%%%%%%%%%%%%%%%%

\newpage
\hypertarget{rng-reg-summ}{}
\section{Register Summary}
\label{sec:rng-reg-summ}

The addresses in this section are relative to Random Number Generator base address provided in Table \ref{tab:sysmem-base-address} in Chapter \ref{mod:sysmem} \textit{\nameref{mod:sysmem}}.

The abbreviations given in Column \textbf{Access} are explained in Section \hyperref[glossary-access-types]{\textit{Access Types for Registers}}.

\subfile{\modulefiles/30-RNG-reg-summ__EN}

%%%%%%%%%%%%%%%%%
%%% Registers %%%
%%%%%%%%%%%%%%%%%

\newpage
\section{Registers}

The addresses in this section are relative to Random Number Generator base address provided in Table \ref{tab:sysmem-base-address} in Chapter \ref{mod:sysmem} \textit{\nameref{mod:sysmem}}.

For how to program reserved fields, please refer to Section \hyperref[programming-reserved-field]{\textit{Programming Reserved Register Field}}.

\subfile{\modulefiles/30-RNG-reg__EN}


\end{document}

